% =====================================================================
% Bagian 2 — Studi Literatur dan Posisi Penelitian  (target 1,25 halaman)
% Tabel 1 WAJIB memuat RFFNet sebagai prior art terdekat, dengan
% diferensiasi moda kegagalan — BUKAN klaim "kami lebih baik".
% =====================================================================

\section{Studi Literatur dan Posisi Penelitian}

\subsection{Deteksi dini karhutla dan fusi RGB--termal}

Kebakaran gambut berbeda fundamental dari kebakaran hutan mineral karena
cenderung membara di bawah permukaan (\emph{smoldering}) dengan nyala rendah dan
asap persisten, sehingga lebih sulit dideteksi maupun dipadamkan
\cite{ballhorn2016}. Karakteristik inilah yang membuat pemantauan berbasis citra
optik saja sering gagal pada fase paling awal. Tolok ukur citra udara yang
menjadi rujukan utama, FLAME \cite{flame2021}, FLAME~2 \cite{flame2_2022}, dan
FLAME~3 \cite{flame3_2025}, seluruhnya berfokus pada kebakaran bernyala di
hutan beriklim sedang dan tidak mencakup gambut tropis; tinjauan komprehensif
juga menegaskan kelangkaan dataset besar dan beragam \cite{review2025}.

Citra termal menangkap gradien suhu dan menembus asap yang membutakan citra
optik, sedangkan RGB menyediakan konteks visual untuk mengenali api dan latar,
sehingga fusi keduanya konsisten mengungguli pendekatan satu kanal
\cite{zhang2025}. Pendekatan kami termasuk kategori \emph{feature-level fusion}
dengan mekanisme \emph{cross-attention} \cite{vaswani2017}, \emph{backbone}
Vision Transformer \cite{dosovitskiy2021}, dan representasi swa-awasan DINOv2
\cite{oquab2024} yang menyediakan representasi kuat dengan kebutuhan label
rendah.

Perlu dibedakan dari pendekatan yang menghasilkan representasi mirip-termal
secara sintetis langsung dari citra RGB \cite{synthetic2026}. Pendekatan itu
berguna ketika sensor termal fisik tidak tersedia, namun tidak menangkap suhu
aktual, indikator paling dini sebelum nyala atau asap tebal tampak. Konsistensi
dengan kritik ini kami pegang sebagai batasan metodologis: kanal termal dalam
penelitian ini tidak pernah disintesis dari RGB, termasuk pada augmentasi
degradasi.

\subsection{Ketahanan terhadap modalitas yang hilang}

Kajian komprehensif pertama mengenai perilaku Transformer terhadap data tidak
lengkap menemukan bahwa model kelas ini sensitif terhadap modalitas hilang, dan
strategi fusi yang berbeda memengaruhi ketahanannya secara signifikan
\cite{ma2022}. \emph{Modality dropout}, membuang modalitas secara acak dengan
probabilitas tertentu selama pelatihan, pertama diusulkan pada pengenalan
gestur multimodal \cite{neverova2016} dan sejak itu terbukti efektif lintas
domain, termasuk dengan \emph{encoder} terlatih yang dibekukan dan modul fusi
ringan \cite{dropout2025}. Dalam konteks karhutla, ketahanan ini menjadi syarat
operasional, bukan fitur tambahan: saat asap pekat atau operasi malam hanya
termal yang andal, sedangkan pada perangkat tanpa sensor termal hanya RGB yang
tersedia.

\subsection{Posisi penelitian dan \emph{research gap}}

Tabel~\ref{tab:related} membandingkan penelitian ini dengan empat karya terdekat.
Perbandingan disusun menurut \emph{moda kegagalan yang ditangani}, bukan menurut
klaim performa, karena karya-karya tersebut dievaluasi pada dataset dan tugas
berbeda sehingga angkanya tidak sebanding secara langsung.

\begin{table}[!htbp]
\caption{Perbandingan dengan karya terdekat dan posisi \emph{research gap}.}
\label{tab:related}
\centering
\small
\begin{tabular}{p{2.4cm}p{3.6cm}p{5.8cm}}
\hline
\textbf{Karya} & \textbf{Domain \& modalitas} & \textbf{Moda kegagalan yang ditangani} \\
\hline
RFFNet \cite{rffnet2026} \newline (\emph{prior art} terdekat) &
Kebakaran hutan, RGB+termal berpasangan, anotasi piksel &
\textbf{Oklusi parsial}: asap atau halangan menutupi sebagian pandangan
sementara \emph{kedua sensor tetap hidup}; ditangani lewat desain arsitektur
fusi \\
Zhang dkk.\ \cite{zhang2025} &
Kebakaran hutan umum (UAV), RGB+termal nyata &
Tidak menguji ketahanan terhadap modalitas hilang; disasar untuk kebakaran
bernyala, bukan rezim membara \\
Transfer learning gambut \cite{peat2026} &
Kebakaran gambut (domain terdekat), RGB saja &
Tanpa sinyal termal; adaptasi domain lewat \emph{fine-tuning}, bukan fusi \\
Termal sintetis \cite{synthetic2026} &
Kebakaran umum, kanal termal disintesis dari RGB &
Tidak menangkap suhu aktual, indikator paling dini yang justru disasar
penelitian ini \\
\hline
\textbf{Penelitian ini} &
Rezim membara sebagai sasaran rancangan, RGB+termal, drone &
\textbf{Kehilangan modalitas total}: satu sensor mati, gagal, atau tidak
terpasang; ditangani lewat \emph{modality dropout} eksplisit dan diukur lewat
kurva degradasi kontinu \\
\hline
\end{tabular}
\end{table}

Diferensiasi terhadap RFFNet perlu dinyatakan dengan tepat karena keduanya
bekerja pada data yang sama, anotasi piksel RFFNet juga menjadi sumber label
penelitian ini. \textbf{Keduanya menangani moda kegagalan yang berbeda dan
saling melengkapi, bukan bersaing.} RFFNet menjawab situasi ketika informasi
tersedia namun tertutup sebagian; penelitian ini menjawab situasi ketika satu
kanal informasi hilang sepenuhnya, dan sistem lapangan yang matang membutuhkan
keduanya. Kami tidak mengklaim keunggulan performa atas RFFNet, dan perbandingan
langsung terhadap bobot RFFNet terlatih tidak terlaksana dalam jendela waktu
kompetisi (Subbagian~4.11).

Irisan yang belum diisi keempat karya di atas adalah kombinasi tiga hal
sekaligus: fusi RGB--termal yang dirancang tahan terhadap kehilangan sensor
total, protokol evaluasi ketahanan yang kontinu alih-alih biner, dan lokalisasi
yang tetap berjalan tanpa label piksel.
